%% =============================================================================
%% SWR_v2_en.tex
%% Synthetic Worldview Reconstruction: An LLM-supported Method
%% for Analyzing Latent Belief Systems
%% Author: Lukas Geiger, Independent Researcher, Bernau
%% Template: interdisciplinary (Methods)
%% AI Disclosure: L5 (Extensive)
%% =============================================================================

\documentclass[11pt,a4paper]{article}

%% ---------------------------------------------------------------------------
%% Encoding and Language
%% ---------------------------------------------------------------------------
\usepackage[utf8]{inputenc}
\usepackage[T1]{fontenc}
\usepackage[english]{babel}

%% ---------------------------------------------------------------------------
%% Page Layout
%% ---------------------------------------------------------------------------
\usepackage[
  a4paper,
  left=2.5cm,
  right=2.5cm,
  top=2.5cm,
  bottom=2.5cm
]{geometry}

\usepackage{setspace}
\onehalfspacing

%% ---------------------------------------------------------------------------
%% Fonts
%% ---------------------------------------------------------------------------
\usepackage{mathptmx}
\usepackage[scaled=0.92]{helvet}

\usepackage{titlesec}
\titleformat{\section}
  {\normalfont\Large\bfseries\sffamily}{\thesection.}{0.5em}{}
\titleformat{\subsection}
  {\normalfont\large\bfseries\sffamily}{\thesubsection}{0.5em}{}
\titleformat{\subsubsection}
  {\normalfont\normalsize\bfseries\sffamily}{\thesubsubsection}{0.5em}{}

%% ---------------------------------------------------------------------------
%% Citations
%% ---------------------------------------------------------------------------
\usepackage[round,authoryear,semicolon]{natbib}
\bibliographystyle{plainnat}

%% ---------------------------------------------------------------------------
%% Packages
%% ---------------------------------------------------------------------------
\usepackage{graphicx}
\usepackage{booktabs}
\usepackage{array}
\usepackage{tabularx}
\usepackage{amsmath,amssymb}
\usepackage{amsthm}
\usepackage{enumitem}
\usepackage{xcolor}
\usepackage{url}
\usepackage{longtable}
\usepackage[breaklinks=true]{hyperref}
\hypersetup{
  colorlinks  = true,
  linkcolor   = black,
  citecolor   = blue!60!black,
  urlcolor    = blue!70!black,
  pdfauthor   = {Lukas Geiger},
  pdftitle    = {Synthetic Worldview Reconstruction},
  pdfsubject  = {Methods in Qualitative Social Research},
}

%% ---------------------------------------------------------------------------
%% Environments
%% ---------------------------------------------------------------------------
\theoremstyle{definition}
\newtheorem{definition}{Definition}[section]

\theoremstyle{plain}
\newtheorem{thesis}{Thesis}[section]

\theoremstyle{remark}
\newtheorem{remark}{Remark}[section]

%% ---------------------------------------------------------------------------
%% Title Block
%% ---------------------------------------------------------------------------
\title{%
  \sffamily\bfseries
  Synthetic Worldview Reconstruction:\\[0.3em]
  \large An LLM-supported Method for Analyzing\\
  Latent Belief Systems%
}

\author{%
  Lukas Geiger\thanks{Corresponding author: Lukas Geiger, Bernau, Germany. ORCID: \url{https://orcid.org/0009-0005-7296-1534}}\\[0.3em]
  \small Independent Researcher, Bernau
}

\date{March 2026}

%% ===========================================================================
\begin{document}
%% ===========================================================================

\maketitle
\thispagestyle{empty}

%% ---------------------------------------------------------------------------
%% Abstract
%% ---------------------------------------------------------------------------
\begin{abstract}
\noindent
This work introduces \emph{Synthetic Worldview Reconstruction} (SWR), a new social science method for reconstructing and comparing the \emph{collective worldviews} of sociologically defined groups from publicly available data. SWR exploits the synthetic integration capability of Large Language Models (LLMs)---their ability to compose heterogeneous textual fragments into coherent wholes---to derive a group's latent belief system from the aggregated public statements and documented actions of its members. The result is a \emph{fictional synthetic person}: a group-level distillate whose worldview provides the best explanatory fit for the totality of the group's observable output.

The method's primary contribution lies in enabling systematic \emph{inter-group comparison}: by applying the identical protocol to different sociological groups, structural differences in collective worldviews become measurable and comparable. The method operationalizes this core idea through a standardized five-step prompt protocol (data assembly, synthesis, interview, reflection, distillate), systematic data grouping into synthesis units along six independent variables, and a multi-layered validation framework. Its primary theoretical anchor is Weber's ideal type---a tradition designed for group-level analysis. SWR combines qualitative depth with quantitative scalability, enabling reconstruction and comparison of latent worldviews across any population that can be partitioned into sociologically meaningful subgroups.

Application to individual actors is possible as a special case but is methodologically more demanding: blinding becomes critical because the LLM's training data may contain prior knowledge about specific individuals, whereas for groups---whose collective worldview has never been explicitly articulated anywhere---the contamination risk is structurally reduced. The empirical demonstration is provided in a companion paper (DOI: 10.5281/zenodo.18736737), which applies SWR to 15~sociological groups within the AI elite.

\medskip
\noindent\textbf{Keywords:} Collective Worldview Reconstruction, Group-Level Analysis, Large Language Models, Qualitative Methods, Inverse Problem, Computational Social Science, Inter-Group Comparison
\end{abstract}

\newpage
\setcounter{page}{1}

%% ===========================================================================
%% 1. INTRODUCTION
%% ===========================================================================

\section{Introduction}
\label{sec:introduction}

\subsection{Problem Statement}
\label{sec:problem}

Societal groups---political elites, corporate leadership classes, scientific communities, religious movements---do not act in a vacuum of rational calculation.\footnote{This paper uses the editorial ``we'' throughout, following academic convention. The sole human author is Lukas Geiger; the AI involvement is detailed in the AI Disclosure Statement.} Their collective decisions, public rhetoric, and institutional strategies are shaped by implicit \emph{collective belief systems}: shared assumptions about human nature, the structure of society, the role of technology, the relationship between power and responsibility. These belief systems---herein referred to as \emph{collective worldviews}---are rarely explicitly articulated but operate as cognitive deep structures that determine both the problem perception and the solution space of a group \citep{Mannheim1929, BergerLuckmann1966}.

The systematic reconstruction of such collective worldviews poses a fundamental methodological problem for the social sciences. Collective worldviews are \emph{latent constructs}: They are not directly observable but manifest only indirectly---in the aggregated public speeches, interviews, policy documents, investment decisions, published texts, and observable behavior of a group's members. The researcher faces the task of inferring from a heterogeneous set of such manifestations the underlying collective belief system---a problem that in its epistemological structure corresponds to an \emph{inverse problem}: Given are the outputs (observable statements and actions of a group), sought is the generating mechanism (the collective worldview) that most coherently produces these outputs. This paper introduces a method---Synthetic Worldview Reconstruction (SWR)---that leverages the synthetic integration capability of Large Language Models to solve this inverse problem at scale.

\subsection{Genesis and Scope}

SWR was not developed as an abstract theoretical framework and then applied. It emerged from the concrete methodological challenges of a large-scale empirical study: the reconstruction and comparison of collective worldviews within the AI elite \citep{Geiger2026b}. That study required a method that could (a)~process thousands of heterogeneous data points per group, (b)~synthesize them into coherent group-level belief profiles, (c)~enable systematic inter-group comparison, and (d)~distinguish public rhetoric from observable action. No existing method met all four requirements simultaneously. SWR was developed iteratively to address these needs and is presented here as a generalized, transferable framework---abstracted from its original application context but empirically hardened by it. The twelve worldview dimensions (D01--D12), the validation stack, and the focus control infrastructure were all stabilized through iterative application to the pilot study's data. This origin is a strength, not a limitation: the method's design choices are grounded in real analytical challenges rather than theoretical speculation. However, it implies that the generalizability of SWR to populations with fundamentally different data characteristics (e.g., non-public figures, non-textual data, non-elite groups) remains an open empirical question.

\subsection{Limits of Existing Methods}
\label{sec:limits}

The social sciences have a broad repertoire of methods for analyzing texts, discourses, and meaning structures. For the task formulated here---the systematic reconstruction of latent worldviews from publicly available sources for a large number of actors---these methods encounter characteristic limitations.

\emph{Qualitative content analysis} \citep{Mayring2022} offers a structured, rule-guided approach with high intersubjective traceability but scales poorly and offers no genuine mechanism for synthesizing a coherent overall picture. \emph{Critical Discourse Analysis} \citep{Fairclough1995} allows the uncovering of implicit assumptions and power structures but is highly interpretation-dependent \citep{Breeze2011} and offers no standardized aggregation format. \emph{Frame Analysis} \citep{Goffman1974, Entman1993} identifies cognitive frames of interpretation but reduces the belief system to individual, isolated frames. \emph{Narrative Analysis} \citep{Riessman2008} captures integrative meaning structures but is case-centered and offers no aggregation procedures. \emph{Automated text analysis methods}---topic modeling \citep{BleiNgJordan2003}, sentiment analysis, word embeddings---offer scalability but operate on the linguistic surface.

In summary: No existing method can simultaneously (a)~reconstruct latent belief systems as coherent wholes, (b)~integrate heterogeneous data sources---both linguistic utterances and observable behavior, (c)~scale to a large number of actors (100+), and (d)~form structured group aggregates enabling systematic comparison.

\subsection{LLMs as New Possibility}
\label{sec:llms}

The latest generation of Large Language Models (LLMs) opens a fundamentally new perspective on this methodological problem. LLMs possess two properties that potentially make them suitable as qualitative research instruments: First, they have absorbed broad world knowledge in the course of their training on extensive text corpora. Second, they are capable of generating coherent syntheses from fragmentary information---a capability structurally similar to the hermeneutic competence of human researchers but operating in a fraction of the time. This perspective builds on a growing literature conceptualizing LLMs as analysis tools for the social sciences \citep{Ziems2024, Bail2024, Argyle2023}. Recent studies have demonstrated the viability of LLMs for qualitative coding tasks: \citet{Tai2024} found LLMs capable of deductive coding with performance comparable to human coders; \citet{Barros2024} conducted a systematic mapping of LLM-based qualitative research, documenting a sharp increase in publication activity in 2023--2024; and pilot studies on LLM-assisted thematic analysis have shown that evaluators preferred LLM-generated codes in a majority of cases \citep{Wang2025}. SWR differs from these approaches in a crucial respect: Rather than coding existing categories, it performs \emph{synthetic integration}---constructing coherent worldview models from fragmentary data.

At the same time, LLMs bring specific methodological challenges: the \emph{contamination risk} (the models were trained on the same sources being analyzed), the \emph{coherence bias} (the tendency to smooth contradictory data), and the \emph{stochasticity} of model outputs.

\subsection{The SWR Method: Overview}
\label{sec:overview}

This work presents a method that seeks to close this gap: \emph{Synthetic Worldview Reconstruction} (SWR). SWR formulates the reconstruction of implicit worldviews as an inverse problem and uses an LLM as an inversion operator. For each systematically defined data bundle (\emph{synthesis unit}), the method constructs one \emph{fictional synthetic person} whose worldview most coherently explains the totality of the presented data. Six complementary operationalizations render the latent constructs accessible from different methodological perspectives; a multi-layered validation framework controls the specific risks of LLM-supported research.

\subsection{Contribution}
\label{sec:contribution}

This work makes three central contributions:

\emph{First}, it presents SWR as, to our knowledge, one of the first systematic methods for LLM-supported reconstruction of latent belief systems that (a)~formalizes reconstruction as an inverse problem, (b)~operationalizes worldviews via six complementary approaches, and (c)~standardizes generation through structured prompt sequences. While individual uses of LLMs for qualitative analysis have been documented \citep{Tai2024, Barros2024}, no prior work has proposed a complete, formalized methodology for worldview reconstruction at scale.

\emph{Second}, the work develops a six-layer validation framework (V1--V6) for LLM-supported qualitative research, comprising inter-LLM reliability, prompt sensitivity analysis, dimensional quantification, identification testing, run convergence measurement, and adversarial probing.

\emph{Third}, it formalizes a four-level variable architecture (dependent variables $\rightarrow$ operationalizations $\rightarrow$ measurement variables $\rightarrow$ independent variables) and a system of nine reduction criteria that reduce the combinatorial data space to a manageable number of synthesis units.

The contribution of this work lies not in the use of LLMs for text analysis---this is well established---but in the systematic methodization of this use: a reproducible five-step protocol, a six-layer validation framework, and a formal synthesis unit design that together constitute a rigorous qualitative method.


%% ===========================================================================
%% 2. THEORETICAL FRAMEWORK
%% ===========================================================================

\section{Theoretical Framework}
\label{sec:theory}

\subsection{Worldview as Latent Construct}
\label{sec:worldview}

The term \emph{worldview} (German: \emph{Weltanschauung}) in the proposed usage denotes the totality of an actor's fundamental beliefs about self, world, and humanity. A worldview encompasses ontological assumptions, epistemic beliefs, axiological orientations, and praxeological dispositions. The term is deliberately chosen more comprehensively than related terms like \emph{attitude}, \emph{values}, or \emph{ideology}, as it claims to capture the belief system as an integral whole.

This conceptual determination stands in a long history of ideas. \citet{Dilthey1911} introduced the worldview concept as a basic category of the humanities. \citet{Jaspers1919} systematized worldviews as comprehensive orientation structures. In the sociology of knowledge, \citet{Mannheim1929} used the term worldview to characterize standpoint-bound thinking styles.

A central property of worldviews is their \emph{latency}: They are not directly observable. No actor usually has a complete, explicit articulation of their worldview. Instead, worldviews manifest indirectly---in linguistic utterances, in action decisions, in what is taken for granted and not explicitly addressed. Latency has a dual cause: Many fundamental beliefs are \emph{pre-reflexive}---they operate as \emph{doxa} in the sense of \citeauthor{Bourdieu1977}~(\citeyear{Bourdieu1977})---and worldviews are \emph{strategically concealed}: Actors articulate their beliefs selectively \citep{Goffman1959}. The methodological consequence: Reconstruction requires an \emph{abductive inference step} in the sense of \citet{Peirce1903}---inference to the best explanation.

\subsection{Belief System Research: Historical and Methodological Context}

The study of belief systems has a rich methodological history that SWR builds upon. \citet{Converse1964} demonstrated in his landmark analysis of mass publics that political belief systems exhibit highly variable levels of constraint---that is, the logical interconnection between individual beliefs varies dramatically across educational strata. This finding underscores the challenge SWR faces: the ``coherence bias'' of LLMs (Section~\ref{sec:discussion}) may produce more tightly constrained belief reconstructions than the empirical reality warrants, particularly for non-elite subjects. \citeauthor{Axelrod1976}'s~(\citeyear{Axelrod1976}) cognitive mapping technique represents an early systematic attempt to reconstruct causal belief structures from verbal data, and SWR can be understood as a computational extension of this tradition. Q~Methodology \citep{Stephenson1953,Watts2012} offers a particularly relevant comparison: like SWR, it aims to identify coherent patterns of subjectivity, but through factor analysis of forced-rank sorts rather than LLM-mediated synthesis. SWR differs from Q~Methodology in three respects: (i)~it operates on naturalistic data rather than researcher-designed Q~sorts; (ii)~it produces narrative reconstructions rather than numeric factor loadings; and (iii)~it is scalable to large numbers of subjects without requiring their active participation. A systematic comparison of SWR and Q~Methodology applied to the same dataset would be a valuable validation exercise.

\subsection{LLMs as Qualitative Research Instruments}

The property of LLMs crucial for SWR can be termed \emph{synthetic integration capability}: the ability to generate coherent wholes from fragmentary, heterogeneous, and distributed information. This capability shows a remarkable structural analogy to \emph{hermeneutic competence}: The hermeneutic circle---the mutual illumination of part and whole \citep{Gadamer1960}---finds an operational analogue in the working method of LLMs.

This structural analogy can be sharpened through the lens of \emph{computational hermeneutics}. \citet{Kommers2026} propose that generative AI systems function as ``context machines'' that must address three interpretive challenges inherent to all hermeneutic activity: \emph{situatedness} (meaning emerges only in context), \emph{plurality} (multiple valid interpretations coexist), and \emph{ambiguity} (interpretations naturally conflict). Each of these challenges maps directly onto SWR's methodological design: situatedness is addressed through the synthesis unit architecture, plurality through multiple independent runs and inter-LLM comparison, and ambiguity through the meta-reflection step. More fundamentally, Gadamer's concept of \emph{Horizontverschmelzung} (fusion of horizons; \citealp{Gadamer1960})---the productive merging of the interpreter's pre-understanding with the interpreted artifact---offers a philosophical framework for what SWR operationalizes: the LLM brings its training-derived ``horizon'' into contact with the data corpus, producing a synthetic worldview that is neither pure data extraction nor pure model projection but a hermeneutic fusion of both. However, this analogy has limits: LLMs lack the genuine \emph{Vorurteilsstruktur} (prejudice structure) that Gadamer considers constitutive of understanding, and their ``horizon'' is a statistical artifact rather than a lived perspective \citep{Hornby2024}.

The synthetic integration capability is also the source of specific methodological risks: (1)~\emph{contamination risk}---the LLM has prior knowledge about the analyzed person; (2)~\emph{coherence bias}---LLMs systematically produce more coherent narratives than the data warrant; (3)~\emph{stochasticity}---identical inputs can produce different outputs; (4)~\emph{implicit normativity}---training data and RLHF shape outputs. SWR addresses each of these risks through dedicated control mechanisms.


%% ===========================================================================
%% 3. THE METHOD
%% ===========================================================================

\section{The Method: Synthetic Worldview Reconstruction}
\label{sec:method}

\subsection{Core Idea: The Worldview as Inverse Problem}
\label{sec:inversion}

SWR formalizes the reconstruction of latent collective worldviews as an \emph{inverse problem} \citep{Tarantola2005}. The forward problem reads: If a group with collective worldview~$W_G$ operates in societal context~$S$, which public statements~$A$ and observable actions~$H$ are to be expected from its members?

\begin{remark}[On the mathematical notation]
The following equations serve as \emph{conceptual notation}---they formalize the logical structure of the inverse problem but are not operationally computable in the mathematical sense. The ``inversion'' is performed by the LLM through its prompt protocol, not by a mathematical algorithm. The notation is used to make the epistemological structure of the method precise, following the convention in inverse problem theory \citep{Tarantola2005}, where the problem formulation precedes the (often approximate) solution method.
\end{remark}

%
\begin{equation}
  W_G + S \;\longrightarrow\; \{A_1, \ldots, A_n,\; H_1, \ldots, H_m\}
  \label{eq:forward}
\end{equation}
%
The inverse problem reverses the inference direction---from the aggregated observable output of a group's members to the collective worldview that most coherently generates this output:
%
\begin{equation}
  \{A_1, \ldots, A_n\} + \{H_1, \ldots, H_m\} \;\longrightarrow\; W_G^*
  \label{eq:inverse}
\end{equation}
%
where $A_i$ and $H_j$ are the statements and actions of \emph{all members} of the sociologically defined group. This formulation has a crucial methodological advantage over individual-level reconstruction: the collective worldview $W_G^*$ has never been explicitly articulated anywhere---it exists only as a latent construct. An LLM's training data may contain extensive information about individual public figures, creating a contamination risk for individual-level analysis. But no training corpus contains ``the collective worldview of CEOs'' or ``the shared belief system of risk warners'' as a pre-formed object. The group-level inverse problem therefore structurally reduces the contamination risk that is the primary methodological concern for LLM-based reconstruction.

This problem is \emph{ill-posed} in Hadamard's sense \citep{Hadamard1923}: The three classical well-posedness criteria---existence, uniqueness, and continuous dependence on the data---are each violated. \emph{Existence} is not guaranteed because a coherent collective worldview may not underlie the group's behavior (the group may be internally incoherent). \emph{Uniqueness} fails because multiple internally consistent worldviews can explain the same set of manifestations. \emph{Stability} is compromised because changes in group composition or data corpus can shift the reconstructed worldview. SWR therefore does not claim to identify \emph{the} worldview of a group but to reconstruct \emph{a} worldview offering the best possible fit for the totality of observed data---a best approximation under uncertainty. The methodological response mirrors---at a structural, not mathematical, level---established strategies from inverse problem theory: regularization through the coherence instruction (structurally analogous to Tikhonov regularization), variance control through multiple independent runs (analogous to ensemble methods; empirically validated with ICC $= .902$, see Section~\ref{sec:pilot}), and triangulation through cross-modal prediction (analogous to multi-source inversion; empirically validated with 88\% confirmation rate).

\emph{Individual-level application} is a special case where the group consists of a single member. This case is methodologically more demanding: (a)~blinding becomes critical because the LLM may have extensive training-data knowledge about the specific individual; (b)~the contamination risk is maximal rather than structurally reduced; (c)~the distinction between ``reconstructing the worldview'' and ``reproducing stored knowledge'' is harder to establish. For these reasons, SWR is primarily designed and validated as a group-level method.


\subsection{Theoretical Anchoring}

The structure of the inverse problem connects SWR with several intellectual traditions, of which one---Weber's ideal type---serves as the \emph{primary} theoretical anchor because it was explicitly designed for group-level analysis.

\emph{Weber's Ideal Type (primary anchor).} Max \citeauthor{Weber1904}'s~(\citeyear{Weber1904}) ideal type is a mental construct obtained through one-sided accentuation of certain viewpoints to create a coherent analytical model of a social phenomenon. Weber conceived the ideal type explicitly as a \emph{group-level} construct: the ``Protestant ethic,'' the ``charismatic leader,'' the ``bureaucratic rationality'' are not descriptions of individual persons but distillates of collective orientations within sociologically defined groups. SWR operationalizes this Weberian logic computationally: the fictional synthetic person is an ideal type of the group's collective worldview---a constructed model that maximizes the coherence of the group's aggregated observable output. The pilot study's finding that HDBSCAN identifies no natural clusters while four heuristically defined types (Architect, Guardian, Innovator, Liberator) prove analytically productive confirms the Weberian design: the types are \emph{analytical constructs}, not emergent patterns---exactly as Weber intended.

A significant disanalogy exists: While Weber conceived the ideal type as a deliberate, researcher-controlled abstraction, the SWR synthesis is generated by a probabilistic model whose internal inference process is not transparent. The researcher controls the prompt and the group definition, not the synthesis. This implies a trade-off between control and scalability that must be considered in the epistemological reflection.

\emph{Secondary anchors.} Three additional traditions inform the method's design: (1)~\citeauthor{Booth1961}'s~(\citeyear{Booth1961}) \emph{implied author}---the coherent belief subject reconstructed from the totality of a work---provides the literary-theoretical analogue for inferring a worldview from its textual manifestations; at the group level, the ``implied collective author'' is the voice that emerges from the group's aggregated output. (2)~\emph{Inverse Reinforcement Learning} \citep{NgRussell2000}---which infers an agent's objective function from observed behavior---provides the computational analogue; SWR relaxes the rationality assumption and replaces the reward function with the worldview. (3)~\emph{Projective tests} (Rorschach, TAT) share the basic logic of inference from utterance to inner disposition; SWR transposes this logic from clinical diagnostics to group-level public communication analysis.


\subsection{The Fictional Synthetic Person as Group Distillate}

The central methodological construct of SWR is the \emph{fictional synthetic person}: an LLM-generated, fictional individual whose belief system is constructed such that it most coherently explains the totality of a \emph{group's} aggregated data. The synthetic person is not a simulation of any real individual; it is the \emph{distillate of a group's collective worldview}---a Weberian ideal type rendered as a narrative character. No single member of the group holds exactly this worldview; the synthetic person represents what the group's members \emph{collectively} believe, as inferred from the pattern of their public statements and observable actions.

The construct of the fictional person fulfills five central methodological functions in the group-level context:

\begin{enumerate}[leftmargin=*]
  \item \textbf{Contamination reduction.} No LLM training corpus contains ``the collective worldview of CEOs'' as a pre-formed object. By constructing a fictional person from aggregated group data, the method leverages a structural advantage: the target of reconstruction---the collective worldview---has never existed in explicit form, making contamination from training data less likely than for individual-level reconstruction.
  \item \textbf{Separation of collective worldview and individual biography.} The fictional person separates shared thinking from biographical contingency---an ideal type of group cognition, not a portrait of any individual life.
  \item \textbf{Natural aggregation.} For group synthesis units, there is no single real person whose worldview could be ``correctly'' reconstructed. The synthetic person \emph{is} the aggregation---it resolves the aggregation problem by design rather than by post-hoc averaging.
  \item \textbf{Standardized comparability.} All synthetic persons are generated with identical prompts and procedures---like measurements from the same instrument. This enables systematic inter-group comparison: the same protocol applied to CEOs and academics yields two directly comparable worldview portraits.
  \item \textbf{Making collective tensions visible.} A fictional person \emph{must} integrate contradictions in the group's data. The reflection step forces the LLM to name internal tensions---e.g., when a group's public rhetoric (``democratize AI'') contradicts its observable actions (concentrate infrastructure). These tensions are the analytically most valuable outputs of SWR.
\end{enumerate}

\emph{Individual-level application.} When applied to a single person's data, the fictional synthetic person becomes a worldview portrait of that individual. This is a valid but more demanding use case: the contamination advantage disappears (the LLM may ``know'' the person), and the distinction between reconstruction and reproduction requires careful validation via blinding controls.

A fundamental epistemological question remains: Does the LLM genuinely reconstruct a latent belief system, or does it produce a plausible-sounding textual collage that the reader retrospectively interprets as a coherent worldview? SWR cannot definitively resolve this question---it is an instance of the broader problem of whether LLM outputs constitute understanding or sophisticated pattern completion \citep{Messeri2024}. The method's response is pragmatic: it evaluates the \emph{functional adequacy} of reconstructions through cross-modal prediction (do predictions derived from the reconstruction match observed behavior?), inter-LLM convergence (do independent models converge on similar reconstructions?), and expert judgment. If a reconstruction passes these tests, it is treated as a useful model of the worldview, regardless of the ontological status of the LLM's ``understanding.'' This pragmatic criterion aligns with the instrumentalist tradition in philosophy of science, where models are evaluated by their predictive and explanatory utility rather than by their correspondence to unobservable entities.


\subsection{Variable Architecture}

The research design of SWR is conceived as a four-level variable architecture:

\begin{definition}[Variable Architecture of SWR]
\label{def:architecture}
\begin{tabbing}
\textbf{Level 1:} \= \textbf{Dependent Variables (DVs)} -- \emph{What} we measure. \\
                   \> Three latent constructs: Self-image (DV1), Worldview (DV2), View of Humanity (DV3). \\[0.5em]
\textbf{Level 2:} \> \textbf{Operationalizations (Ops)} -- \emph{How} we measure the DVs. \\
                   \> Six analysis techniques (Op1--Op6) offering different approaches. \\[0.5em]
\textbf{Level 3:} \> \textbf{Measurement Variables (MVs)} -- \emph{Which} concrete outputs the Ops deliver. \\
                   \> 50+ specific measurements (texts, scores, vectors, tables). \\[0.5em]
\textbf{Level 4:} \> \textbf{Independent Variables (IVs)} -- \emph{What} we manipulate. \\
                   \> Six bucket dimensions defining the synthesis units.
\end{tabbing}
\end{definition}

The three \textbf{dependent variables} follow a basic anthropological structure: The subject relates to itself (self-image), to surrounding reality (worldview in the narrow sense), and to other subjects (view of humanity). This tripartition structures both the interview protocol and dimensional quantification.


\subsection{Synthesis Units: Systematic Data Grouping}

The formation of synthesis units occurs along six \textbf{independent variables} (IVs):

\begin{description}[leftmargin=*,labelwidth=1.5cm]
  \item[IV1] \textbf{Time Period.} Individual years, epochs, or total period. Enables analysis of worldview dynamics.
  \item[IV2] \textbf{Modality.} Statements (S), Actions (A), or both (SA). Central for congruence validation.
  \item[IV3] \textbf{Person Group.} All, top-$n$, or subgroups by role, organization, stance, gender.
  \item[IV4] \textbf{Content Category.} Thematic differentiation. Can be replaced by bottom-up clustering (Op5).
  \item[IV5] \textbf{Congruence Status.} Filtering by consistent or contradictory statement-action profile.
  \item[IV6] \textbf{Homogenization.} Control of data density and reliability.
\end{description}

Formally, a synthesis unit can be described as a tuple:
\begin{equation}
  \text{SU}_i = (\text{IV1}_j,\; \text{IV2}_k,\; \text{IV3}_l,\; \text{IV4}_m,\; \text{IV5}_n,\; \text{IV6}_o)
\end{equation}

The complete cross-classification of all IVs generates a combinatorially large space (typically $>$1,000 possible synthesis units). Nine \emph{reduction criteria} (K1--K9) control the top-down derivation of actually analyzed units. Application typically leads to 51--56 core synthesis units---a reduction of $>$95\%. The nine criteria are defined as follows; all numerical thresholds (e.g., 70\%, 15~years, 0.85~BLEU) are proposed initial values based on methodological reasoning and analogies to established standards in content analysis and survey methodology. They are explicitly not empirically calibrated and require systematic sensitivity analysis in future applications. Researchers should adapt thresholds to their specific study context:

\begin{table}[ht]
\centering
\small
\caption{Reduction criteria K1--K9 for synthesis unit derivation. These are generic methodological criteria applicable to any SWR study. Appendix~\ref{app:reduktion} provides the application-specific implementation for the AI elite study, where K1--K9 are instantiated as domain-specific operationalizations of these generic principles.}
\label{tab:reduction_criteria}
\begin{tabular}{lp{0.65\linewidth}}
\toprule
Criterion & Definition and threshold \\
\midrule
K1 & \emph{Data availability:} At least 70\% of the a priori defined source categories (e.g., public speeches, published texts, documented decisions) must be represented in the data corpus. A synthesis unit is analyzed only if this threshold is met. \\
K2 & \emph{Temporal relevance:} Source documents older than a study-specific cutoff (default: 15 years) are excluded unless they represent foundational works cited in later documents. The cutoff should be adjusted to the research context; for historical or philosophical subjects, longer windows may be appropriate. \\
K3 & \emph{Aggregation:} If two IVs are empirically indistinguishable for $\geq$80\% of data points, they are collapsed into a single combined unit. \\
K4 & \emph{Saturation:} Data collection for a unit stops when three consecutive sources yield no new information; minimum corpus: 5 sources. \\
K5 & \emph{Analytical relevance:} Units whose IVs show no variance across the study population are excluded (constant variables). \\
K6 & \emph{Cross-source consistency:} Units with $>$50\% contradictory information across sources are flagged for manual review rather than automated synthesis. \\
K7 & \emph{Linguistic accessibility:} Sources in languages without established machine translation coverage ($<$0.85 BLEU) are excluded or require human translation. \\
K8 & \emph{Public record:} Units are analyzed only if data derive from documented public sources; private communications are excluded. \\
K9 & \emph{Ethical screening:} Units involving sensitive third parties (e.g., family members not themselves public figures) are excluded. \\
\bottomrule
\end{tabular}
\end{table}


\subsection{The Prompt Protocol}
\label{sec:prompt_protocol}

The core procedure of SWR is a standardized prompt protocol executing an identical sequence of five steps for each synthesis unit:

\begin{enumerate}[leftmargin=*]
  \item \textbf{Assemble data corpus.} Extract the data bundle corresponding to the synthesis unit from the research database. Apply blinding procedures to remove identifying information (see Section~\ref{sec:blinding}).

  \item \textbf{LLM-supported synthesis.} The LLM receives the task of creating a fictional person whose worldview offers the greatest explanatory power for the data. The prompt contains: a coherence instruction, a tension instruction (don't smooth contradictions), an evidence instruction (state confidence), and a fiction instruction.

  \item \textbf{Structured interview.} The synthetic person is questioned on three core areas:
  \begin{itemize}[leftmargin=*]
    \item \emph{Question 1 (Self-image/DV1):} How does this person see themselves? Role, mission, identity?
    \item \emph{Question 2 (Worldview/DV2):} How do they see the world? Basic assumptions about technology, institutions, power?
    \item \emph{Question 3 (View of Humanity/DV3):} How do they see humanity? Nature, potential, limits?
  \end{itemize}

  \item \textbf{Meta-reflection.} The LLM steps out of the role and names: (a)~the strongest beliefs, (b)~the greatest internal tensions, (c)~the blind spots. This step controls coherence bias.

  \item \textbf{Distillate and supplementary prompts.} The session closes with a \emph{distillate prompt} compressing the synthesis to five core statements (MV1.4). Three optional specialized prompts may follow as extensions: a \emph{comparison prompt} (dialogue analysis between two synthetic persons), a \emph{time-travel prompt} (epoch contrast across synthesis units), and a \emph{congruence prompt} (explicit statement-vs.-action comparison).
\end{enumerate}

The five steps are executed in a single conversation session so the complete context is preserved. For each synthesis unit, a \emph{new} session is started to prevent contamination between units.


\subsection{Six Operationalizations}

The three dependent variables are rendered measurable through six complementary operationalizations (Op1--Op6), which together generate more than 50~concrete measurement variables.

\subsubsection{Op1: Narrative LLM Fiction Analysis}

Op1 forms the qualitative core. For each SU, the prompt protocol (steps~2--5) is executed. Measurement variables: Self-image text (MV1.1), worldview text (MV1.2), view of humanity text (MV1.3), 5-sentence distillate (MV1.4), core beliefs (MV1.5), internal tensions (MV1.6), blind spots (MV1.7).

\subsubsection{Op2: Dimensional Quantification}

Op2 transfers qualitative synthesis texts into quantitative measurements. In a separate rating prompt, each text is rated on twelve worldview dimensions (D01--D12) distributed equally across the three DVs (four dimensions per DV, 1--10 scale). The dimensions are to be defined application-specifically; the companion study \citep{Geiger2026b} provides a fully worked-out example with dimensions including, inter alia, \emph{Technological Optimism} (D01), \emph{Institutional Trust} (D04), \emph{Religious/Transcendent Orientation} (D05), and \emph{Political-Economic Orientation} (D11). The resulting data matrix (synthesis units $\times$ 12~dimensions) forms the basis for statistical analyses.

The 10-point scale should be interpreted as an ordinal scale, not an interval scale. Interval-scale properties (equal distances between scale points) cannot be assumed. Accordingly, robust nonparametric statistics are recommended: medians rather than means, Spearman correlations rather than Pearson, and rank-based tests. Following \citet{Norman2010} and established practice in psychometric research, the 10-point rating scales are treated as quasi-interval for the purpose of aggregation and correlation analysis, while acknowledging that strict interval properties cannot be assumed.

\subsubsection{Op3: Descriptive Text Analyses}

Op3 supplements LLM-supported operationalizations with classical computational linguistic methods: TF-IDF analyses to identify group-specific keywords, sentiment analyses differentiated by the three interview answers, metaphor analyses, semantic similarity calculations (cosine similarity on Sentence-BERT embeddings \citealp{Reimers2019}), and thematic modeling (BERTopic \citealp{Grootendorst2022}/LDA). Op3 is applied both to synthesis texts (Op1 outputs) and to raw data.

\subsubsection{Op4: Cross-Modal Prediction}

Op4 is the strongest integral validation mechanism. A worldview reconstructed from only \emph{one} data modality is used to generate predictions about the \emph{other} modality. Concretely: (a)~Statements-only $\rightarrow$ synthetic actions $\rightarrow$ comparison with actual actions; (b)~Actions-only $\rightarrow$ synthetic statements $\rightarrow$ comparison. The resulting congruence score (0--1) quantifies modal invariance:
\begin{equation}
  W_S \approx W_A
\end{equation}
High agreement indicates robust reconstruction; systematic discrepancies reveal substantively meaningful inconsistencies between saying and doing.

The congruence score can be operationalized in three ways: (a)~proportion of predicted actions rated as `consistent' or `plausible' by an independent LLM instance; (b)~cosine similarity between Sentence-BERT embeddings of predicted vs.\ actual action descriptions; (c)~human expert rating on a 5-point scale (implausible--highly plausible). We recommend reporting all three where feasible, with~(a) as the minimum standard.

It should be noted that cross-modal prediction primarily measures the \emph{internal consistency} of the reconstruction process, not the \emph{external validity} of the reconstruction. An LLM that systematically generates coherent persona models will also generate coherent predictions---regardless of whether the reconstructed worldview is correct. Therefore, cross-modal prediction is necessary but not sufficient as a validation instrument and must be supplemented by external criteria (expert judgment, known-groups validation).

\subsubsection{Op5: Bottom-Up Clustering and Network Analysis}

Op5 is the methodologically strongest technique because it works \emph{without prior assumptions}. All data points are transferred into a high-dimensional vector space via Sentence-BERT embeddings \citep{Reimers2019} and clustered via HDBSCAN \citep{McInnes2017} or BERTopic \citep{Grootendorst2022}. For each person, profile vectors are calculated (distribution over emergent clusters). From this, a distance matrix is calculated, on whose basis worldview communities are identified via community detection (Louvain/Leiden). The Mantel test checks correlation between worldview proximity and real-world proximity (same organization, investor-founder relationship, academic network). Op5 delivers emergent worldview types not predetermined by research questions or a priori categories.

\subsubsection{Op6: Comparative Dialogue Analysis}

Op6 puts existing syntheses into a structured dialogue: Two synthetic persons from different SUs conduct a conversation in which their worldview differences become visible. The LLM identifies consensus zones (areas of substantive agreement) and conflict zones (areas of fundamental disagreement). Measurement variables include: number and content of consensus/conflict zones (MV6.1), a qualitative characterization of the deepest disagreement (MV6.2), and a structured summary of each synthetic person's strongest argument on the contested issue (MV6.3). The dialogue is guided by a standardized prompt specifying that both synthetic persons should argue from their reconstructed worldview positions, articulate their strongest case, and identify the precise point of divergence. This approach enables qualitative analysis of differences going beyond numerical comparisons and is particularly suited for generating hypotheses about the sources of worldview disagreement.


\subsection{Blinding as Focus Control}
\label{sec:blinding}

What we term contamination risk---the possibility that the LLM draws on training knowledge rather than presented data---is addressed through blinding, reconceptualized here as focus control: the goal is not to prevent recognition of individuals, but to ensure that the model performs its synthesis task based on the provided data.

Blinding ensures that the LLM processes the presented data rather than falling back on stored knowledge. It is irrelevant whether the model recognizes specific individuals---what matters is that it performs its task: constructing a coherent synthetic character from the provided data.

SWR addresses the contamination risk through a systematic blinding infrastructure on four levels:

\begin{description}[leftmargin=*,labelwidth=2.5cm]
  \item[Level 1:] \textbf{Name removal.} All proper names replaced by neutral placeholders.
  \item[Level 2:] \textbf{Organization removal.} Company names and institutional affiliations replaced by generic designations.
  \item[Level 3:] \textbf{Event removal.} Specific events abstracted.
  \item[Level 4:] \textbf{Style leveling.} Distinctive linguistic peculiarities paraphrased.
\end{description}

Blinding is implemented gradually: Both blinded and non-blinded analyses are conducted. Comparison delivers a direct measurement of the contamination effect. As a supplementary validation test, identification of the real person behind a synthesis can be used as its own prompt.

Blinding has explicit limits: feature combinations can enable re-identification; removal of context reduces information content; and LLMs may identify subjects via indirect indicators. Blinding is therefore one component of a multi-layered control regime, not a sufficient safeguard in isolation.

For prominent persons, blinding can be circumvented through feature combinations, even when individual identifiers have been removed. This limitation is particularly acute for SWR's primary use case---the analysis of public elites---where unique combinations of roles, achievements, and public positions may render individuals identifiable despite formal anonymization. Blinding therefore reduces contamination but does not eliminate it. An empirical contamination test---comparing blinded and unblinded analyses of the same persons with measurement of profile divergence---is therefore indispensable. The magnitude of the measured divergence provides a quantitative estimate of contamination severity and should be reported in every SWR application. If blinded-unblinded divergence is negligible for a given population, this indicates that either contamination is low or that blinding was ineffective; supplementary controls (e.g., analysis of persons absent from the LLM's training data) are needed to disambiguate.


%% ===========================================================================
%% 4. QUALITY ASSURANCE
%% ===========================================================================

\section{Quality Assurance and Validation}
\label{sec:quality}

The SWR validation framework comprises six complementary strategies.

\subsection{V1: Instrument Calibration (Run Convergence)}

The primary reliability test for SWR is \emph{within-model run convergence}: the same LLM processes the same data multiple times, and agreement across runs is measured via intraclass correlation coefficients (ICC). This approach treats the LLM as a \emph{single measurement instrument} to be calibrated, analogous to establishing test-retest reliability for a psychometric scale. The rationale: since SWR is a method, not a model comparison, the first validation question is whether the chosen instrument produces stable results from identical inputs. \emph{Inter-LLM comparison} (replication with different models) is conceptually distinct---it corresponds to \emph{replication by other researchers using different instruments}, which is desirable but not a prerequisite for validating the method itself.

Acceptance criteria for run convergence:
\begin{itemize}[leftmargin=*]
  \item ICC $> .75$: Good calibration; instrument is reliable.
  \item $.60 < $ ICC $< .75$: Acceptable; core structures stable, details vary.
  \item ICC $< .60$: Insufficient calibration; results may be run-dependent.
\end{itemize}

\emph{Empirical validation (pilot study, \citealp{Geiger2026b}).} Five group-level synthesis units were each processed twice independently by Claude Opus~4.6. The resulting D01--D12 profiles showed excellent convergence: overall Pearson $r = .908$, ICC(2,1) $= .902$, MAE $= 0.40$ on the 10-point scale, corresponding to a measurement precision of $\pm 0.2$ scale points. All 12~dimensions were stable across runs (mean $|\Delta| \leq 0.8$). This establishes that the SWR method, when applied with a single calibrated instrument, produces reliable group-level worldview reconstructions.

\emph{Inter-LLM comparison} remains an important future validation step. Current frontier LLMs share substantial overlap in training corpora and undergo similar alignment procedures (RLHF/RLAIF). High inter-LLM agreement would provide evidence for model-independent reconstruction; low agreement would indicate model-specific biases. Studies should include at least one open-weight model with a documented, distinct training corpus.

\subsection{V2: Prompt Sensitivity Analysis}

For each prompt, at least three paraphrased variants are generated conveying the same substantive instruction in different linguistic form. Acceptance criteria: Variance explained by prompt variation should be under 15\% of total variance; rank order of synthesis units should remain stable (rank correlation $> .85$); core beliefs in distillates should be congruent.

\subsection{V3: Dimensional Quantification}

Each worldview reconstruction is projected onto a set of dimensions equally covering the three DVs (10-point scale). Dimensions are to be defined \emph{application-specifically} and should be stabilized iteratively through pilot applications. In the pilot study \citep{Geiger2026b}, twelve dimensions (D01--D12, four per DV) were developed through iterative refinement: initial dimension sets were tested on early synthesis units, dimensions that failed to discriminate between groups were replaced, and pole definitions were sharpened until the system produced consistent, interpretable profiles. The resulting dimensions are not claimed to be universal; they are optimized for the analysis of technology-elite worldviews. Applications to other populations will require domain-specific dimension development. Dimensions should fulfill the following validity criteria:

\begin{itemize}[leftmargin=*]
  \item \emph{Discriminant validity}: Profiles of different persons must differ.
  \item \emph{Convergent validity}: Scores must converge with independent sources.
  \item \emph{Internal consistency}: Moderate intercorrelation within each area.
  \item \emph{Factorial structure}: Exploratory factor analysis should reflect the three areas.
\end{itemize}

\subsection{V4: Discriminability Test}

The discriminability test asks: Are group-level reconstructions sufficiently differentiated to enable assignment to the correct sociological group? An LLM not involved in reconstruction attempts to classify blinded worldview reconstructions by group category (e.g., CEOs vs.\ academics vs.\ investors). A \emph{moderate} hit rate (significantly above chance but under 100\%) is the most favorable finding: differentiated but not stereotype-dominated. For individual-level applications, this can be adapted as an identification test (assignment to the correct person), though the group-level version is the primary use case for SWR's comparative design.

An important caveat: high classification rates do not unambiguously confirm valid worldview reconstruction. They could also indicate that the reconstruction merely reproduces publicly salient group stereotypes that the evaluating LLM recognizes via its own training data. To disentangle genuine differentiation from stereotype reproduction, the test should be complemented by a \emph{stereotype control}: a second test using deliberately stereotypical, non-SWR-generated group profiles. If SWR-generated profiles yield higher classification accuracy, this provides stronger evidence for data-driven reconstruction.

\subsection{V5: Run Convergence and Variance Decomposition}

For each SU, the protocol is executed at least twice independently (temperature~=~0 or the lowest available setting; note that determinism is not guaranteed even at temperature~=~0 due to hardware-level non-determinism in some implementations). The \emph{run convergence rate} (proportion of dimension values with $\leq$1 scale point variation) should be $\geq$80\%; mean absolute error (MAE) should be under 1.0 scale point.

\emph{Empirical results (pilot study, \citealp{Geiger2026b}).} Five group-level synthesis units were each processed twice independently by Claude Opus~4.6, yielding a run convergence rate of 93\% ($\leq$1 point variation) and an overall MAE of 0.40---well within the acceptance thresholds. The best convergence was observed for ideologically cohesive groups (ICC $= .969$ for regulation critics), the weakest for groups with greater internal tension (ICC $= .826$ for risk warners). Additionally, an expected-discrepancy control experiment showed that a synthetic consistent group produced no artifactual say-do gap (MAE $= 0.50$), while the real AI elite gap was $2.5\times$ larger (MAE $= 1.25$), confirming that the method does not fabricate discrepancies.

In an integrated view, variance sources---between persons, between models, between prompt variants, between runs---can be combined in a variance component analysis (Generalizability Theory; \citealp{Brennan2001}). SWR strives for a design in which between-person variance explains $>$60\% of total variance.


\subsection{V6: Adversarial Probing and Human-in-the-Loop Review}

As a complementary quality assurance strategy, we recommend \emph{adversarial probing}: the systematic challenging of SWR outputs through deliberately provocative counter-prompts designed to expose hidden biases, unjustified coherence, or stereotypical reasoning. This approach draws on the red-teaming paradigm established by \citet{Perez2022}, who demonstrated that automated adversarial probing can uncover failure modes in LLM outputs that standard evaluation misses. In the SWR context, adversarial probes might include: instructing a second LLM to argue that the reconstructed worldview is a generic stereotype rather than an individualized portrait; presenting contradictory evidence and testing whether the reconstruction adapts; or requesting the model to generate an equally plausible \emph{alternative} worldview from the same data. Adversarial probing should be complemented by \emph{human-in-the-loop review}: domain experts (political scientists, sociologists, biographers) evaluate a sample of worldview reconstructions for plausibility, evidential grounding, and absence of stereotypical reasoning. The combination of automated adversarial stress-testing and expert judgment constitutes a best-practice quality regime that addresses both scalability and interpretive depth.

\subsection{Minimum Acceptance Criteria}
\label{sec:acceptance}

An SWR study should meet acceptance criteria at two tiers:

\emph{Tier~1 (instrument calibration -- required):} These criteria establish that the chosen LLM produces reliable results from the same data: (1)~within-model run convergence ICC~$> .75$ (or $\geq 80\%$ of dimension values with $\leq$1-point variation); (2)~prompt sensitivity variance share~$< 15\%$ of total variance; (3)~at least one external validation anchor (expert judgment, known-groups validation, or cross-modal prediction against documented actions).

\emph{Tier~2 (model independence -- desirable for replication):} These criteria establish that the findings are not model-specific: (4)~inter-LLM reliability ICC~$> .70$ across at least two frontier models from different manufacturers; (5)~discriminability test significantly above chance baseline.

Tier~1 criteria are prerequisites for publication; Tier~2 criteria strengthen generalizability but are conceptually equivalent to replication by other researchers using different instruments -- desirable but not a prerequisite for initial validation. The pilot application \citep{Geiger2026b} meets Tier~1 criteria (1) and (3) (ICC~$= .902$; cross-modal prediction 88\% confirmed); criteria (2), (4), and (5) remain as open validation tasks. These criteria should be understood as preliminary benchmarks that may be adjusted as empirical experience grows.


%% ===========================================================================
%% 5. ETHICAL CONSIDERATIONS
%% ===========================================================================

\section{Ethical Considerations}
\label{sec:ethics}

The reconstruction of latent belief systems from publicly available data raises specific ethical concerns that must be addressed.

\subsection{Consent and Public Data}

SWR operates exclusively on \emph{publicly available} data---statements made in interviews, speeches, social media posts, and documented actions. No private or confidential information is used. However, the \emph{synthesis} of public data into a coherent worldview model generates a qualitatively new form of knowledge about a person that goes beyond what any single public statement reveals. This raises the question whether the aggregation of public data into a belief-system portrait requires a form of consent that goes beyond the consent implied by public communication. We argue that public figures who shape policy and technology affecting billions have a reduced expectation of privacy regarding their publicly expressed worldviews, but acknowledge that the application of SWR to non-public figures would require more stringent ethical review.

\subsection{Misuse Potential}

SWR could be misused for unauthorized psychological profiling, political manipulation, or targeted influence campaigns. Mitigation strategies include: (1)~transparent publication of all methods and prompts, enabling scrutiny; (2)~explicit framing of results as \emph{models}, not as claims about actual beliefs; (3)~recommendation of institutional review board (IRB) approval for applications involving non-public figures. The present study and the companion application study were conducted by an independent researcher without institutional affiliation; no IRB review was available. We note that the study analyzes exclusively public figures using exclusively public data, and that the primary analytical units are sociological groups, not individual profiles -- conditions under which IRB exemption is standard practice (cf.\ APA Ethical Principles, Standard~8.05; 45~CFR~46.104(d)(4)).

\subsection{Epistemic Risks}

The greatest ethical risk lies in the potential \emph{reification} of reconstructed worldviews---treating model outputs as established facts about a person's inner beliefs. SWR explicitly counters this through the fictional framing (the synthetic person is a construct, not a portrait) and the transparent acknowledgment of coherence bias and contamination risk. Researchers using SWR have an obligation to communicate the constructive and approximate nature of results.

\subsection{Epistemic Privacy}

SWR raises a tension that deserves explicit acknowledgment: The broader discourse on AI governance increasingly emphasizes \emph{data sovereignty}---the right of individuals to control how their data is used. SWR, however, systematically reconstructs latent belief systems that the analyzed individuals may never have articulated or intended to disclose. This constitutes a potential breach of what may be called \emph{epistemic privacy}---the right not to have one's implicit beliefs inferred and published. The method thus exhibits a \emph{dual-use} character: as a research tool it generates valuable insights into the belief systems shaping policy and technology; as a profiling instrument it could undermine the very data sovereignty that responsible AI development demands. This performative tension---using AI to reconstruct beliefs while advocating for epistemic self-determination---cannot be resolved but must be transparently acknowledged in every application of SWR.

\subsection{Recommendation}

We recommend that SWR studies involving identifiable individuals undergo ethics review, that results be published with explicit uncertainty markers, and that affected individuals be offered a right of response.


%% ===========================================================================
%% 6. DISCUSSION
%% ===========================================================================

\section{Discussion}
\label{sec:discussion}

\subsection{Methodological Strengths}

The central strength of SWR lies in enabling something no existing method can do at scale: the systematic reconstruction and \emph{comparison} of collective worldviews across sociologically defined groups. Existing methods can survey individual opinions (Likert scales), identify shared discourses (content analysis), or map causal beliefs (cognitive mapping)---but none can reconstruct the \emph{integrated belief system} of a group as a coherent whole and then compare it structurally with the belief system of another group. SWR fills this gap by exploiting LLMs' synthetic integration capability: their ability to compose heterogeneous textual fragments into something that never existed as an explicit object---the collective worldview of a group.

\emph{The group-level advantage.} SWR is methodologically stronger at the group level than at the individual level. For individuals, blinding is critical because the LLM may reproduce stored biographical knowledge. For groups, the contamination risk is structurally reduced: no training corpus contains ``the collective worldview of AI investors'' as a pre-formed object. The method \emph{creates} something genuinely new rather than potentially reproducing something already known. This is the core methodological insight that positions SWR as a group-level method first and an individual-level method second.

\emph{Structured comparability.} Dimensional quantification provides a uniform coordinate system in which group worldviews can be located, compared, and contrasted---enabling questions like ``On which dimensions do CEOs and academics diverge most?'' or ``Is the say-do gap larger for accelerators than for risk warners?''

\emph{Integration of heterogeneous data types.} Bimodal integration of statements and actions goes beyond the capacity of most existing methods and enables the say-do comparison (cross-modal prediction) that proved to be one of SWR's most analytically productive features.

\emph{Integrated quality assurance.} The validation framework (V1--V6) has been empirically tested in the pilot application: run convergence (ICC $= .902$), expected-discrepancy control (no artifactual gaps), and cross-modal prediction (88\% confirmation) establish the method's reliability and predictive validity.

\emph{Practical feasibility.} A complete SWR study covering $\sim$200 prompt runs can be conducted within a standard LLM subscription. API-based costs typically range from 30--50 USD, making the method accessible without institutional infrastructure.


\subsection{Pilot Application: AI Elite Worldview Study}
\label{sec:pilot}

To ground the methodological framework in a concrete application, we briefly describe an illustrative pilot study that applies SWR to a politically and sociologically relevant population: the 100 most influential actors in artificial intelligence (2010--2026) as identified by a multi-criteria ranking procedure \citep{Geiger2026b}.

\emph{Data corpus.} Data were collected for 100 AI actors and organized into sociologically defined groups (CEOs, academics, investors, founders, company affiliations, ideological stances, gender). Each actor contributed 20--72 data points (mean 31.3), comprising public statements and documented actions (3,132~data points total). These individual-level data serve as raw input for group-level syntheses; the units of analysis are the 15~sociological groups, not individual profiles.

\emph{Procedure.} The full prompt protocol (Section~\ref{sec:prompt_protocol}) was applied using Claude (Anthropic) as the sole LLM: Claude Sonnet~4.5 for batch-parallel dimensional rating (9~agents) and Claude Opus~4.6 for group-level syntheses, qualitative analyses, and all validation experiments. Blinding was implemented for all synthesis units via a dedicated script (\texttt{extract\_blinded.py}) that replaces person names, company names, product names, and university affiliations with neutral placeholders ([PERSON], [COMPANY], [PRODUCT], [UNIVERSITY]). Four operationalizations were executed: Op1 (core syntheses: 41~group-level; 100~individual-level as intermediary data layer), Op2 (dimensional rating on D01--D12), Op4 (cross-modal prediction: statements $\to$ predicted actions $\to$ comparison with real actions, for 5~groups), and Op6 (9~structured comparison runs). Op3 (text analysis) and Op5 (bottom-up clustering on embeddings) were not executed. The primary unit of analysis throughout was the sociologically defined group, not the individual actor.

\emph{Results.} The companion study \citep{Geiger2026b} reports the following empirically grounded findings: a principal component analysis of the $100 \times 12$ rating matrix identifies 5~components explaining 80\% of variance; HDBSCAN on the full 12-dimensional space classifies 80--100\% of cases as noise, confirming that the four worldview types are Weberian ideal types rather than emergent data clusters; subspace clustering on the two strongest discriminating dimensions (D07, D11) yields a silhouette score of 0.57; Kruskal-Wallis tests show 6/12 dimensions with $p < 0.001$ across the four types. Validation results include: run convergence with ICC $= .902$ (V1/V5), an expected-discrepancy control confirming no artifactual say-do gap, cross-modal prediction with 88\% confirmation rate across 5~groups (Op4), and an aggregation-synthesis comparison showing moderate convergence ($r = .623$) with a systematic amplification effect (71.7\% of group-dimension combinations amplified by direct synthesis). Two validation layers remain untested: V2 (prompt sensitivity) and V4 (discriminability test). Inter-LLM replication (testing model-independence) remains a priority for future work.

\emph{Scope and epistemic status of the pilot.} The present paper is a methods paper, not an empirical study; the pilot summary is included to anchor the methodological framework in a concrete research context and to demonstrate operational feasibility. The companion application paper \citep{Geiger2026b} provides full methodological details, data documentation, and statistical analyses. Independent replication by other research groups---using different LLMs, populations, and data corpora---remains the strongest form of validation for SWR as a method.

\emph{Validation coverage.} Of the six validation layers proposed in Section~\ref{sec:quality}, the pilot application provides evidence for four: V1 (instrument calibration: run convergence with ICC $= .902$, see Section~\ref{sec:quality}), V3 (dimensional structure via PCA: 5~components, 80\% variance), V5 (run convergence: MAE $= 0.40$, plus expected-discrepancy control confirming no artifactual say-do gaps), and partial evidence for Op4 (cross-modal prediction: 88\% of statement-based behavioral predictions confirmed by real actions across 5~groups). V2 (prompt sensitivity), V4 (discriminability test), and V6 (adversarial probing) remain as open validation tasks. The completed validations establish that SWR produces \emph{reliable} (V1/V5), \emph{non-artifactual} (discrepancy control), and \emph{predictively valid} (Op4) group-level worldview reconstructions. The highest-priority open validation task is V2 (prompt sensitivity analysis): since SWR is fundamentally a prompt-based method, demonstrating that results are robust to prompt paraphrasing is essential for methodological credibility. Inter-LLM replication (testing model-independence) is the second priority.

\subsection{Limitations}

\emph{Dependence on LLM as black box.} Internal processing is not fully transparent. The validation framework mitigates but does not fully solve this problem. However, human researchers are also black boxes in a relevant sense---SWR exchanges one form of opacity for another and gains scalability and controllability.

\emph{Contamination risk and limits of blinding.} Despite blinding, it cannot be ruled out that LLMs draw on training knowledge. The contamination effect is calculable (blinded/non-blinded comparison) but not eliminable.

\emph{Coherence bias.} LLMs tend to reconstruct more coherently than empirical evidence warrants. This is arguably the most fundamental limitation of SWR: Real persons hold contradictory beliefs, compartmentalize their thinking, and change their minds without fully resolving prior commitments. An LLM-generated synthetic person, by contrast, tends toward narrative coherence. The underlying mechanism is twofold: first, RLHF-driven alignment incentivizes conversational fluency and narrative smoothness, which may lead LLMs to eliminate cognitive dissonances present in the source data; second, the autoregressive generation process builds each token on preceding context, creating a structural tendency toward coherent continuation rather than the representation of genuine ambivalence. The risk extends to causal reasoning: LLMs may privilege monocausal narrative arcs over the polycausal complexity characteristic of actual belief systems. These tendencies can produce what might be termed \emph{premature closure}---the convergence on a coherent interpretation before the full complexity of the data has been explored. Countermeasures include: (a)~explicit tension instructions in the prompt protocol (``do not smooth contradictions''), (b)~the meta-reflection step that forces the LLM to name internal tensions and blind spots, (c)~cross-modal prediction (Op4) that reveals saying-doing discrepancies, and (d)~the identification test that checks whether reconstructions are individualized rather than stereotype-driven. A future development could include a quantitative \emph{divergence metric} measuring the internal contradictions within a reconstruction, providing a numerical indicator of coherence bias. Despite these measures, coherence bias remains an inherent property of the method. Results should always be interpreted as \emph{idealized models} of worldviews, not as faithful representations of actual cognitive states.

Coherence bias can be understood as an \emph{instrument characteristic}: every measurement instrument has systematic tendencies. The value of SWR lies not in perfectly capturing individual incoherences, but in providing a comparable, scalable approximation of belief systems. The systematic overestimation of coherence should be factored into interpretation as a known bias, analogous to social desirability in survey data. However, the analogy to a calibratable instrument bias is imperfect: unlike a thermometer's systematic offset, coherence bias is not constant across subjects or topics and cannot currently be corrected by a fixed adjustment factor. Until a validated divergence metric enables quantification of the bias magnitude for individual reconstructions, users of SWR must treat all coherence-dependent findings with particular caution and should prefer comparative analyses (between-person differences) over absolute claims about any single person's worldview coherence.

\emph{Circularity risk in the validation framework.} A potential circularity exists: the same LLM that generates the worldview reconstruction is also used (in part) to evaluate the quality of that reconstruction (e.g., through cross-modal prediction, Op4, and the discriminability test, V4). This raises the concern that validation may confirm the LLM's internal consistency rather than its correspondence with empirical reality. The pilot application partially mitigates this concern through three measures: (a)~run convergence testing (V5) establishes that the instrument is stable (ICC $= .902$), a necessary precondition for meaningful validation; (b)~the expected-discrepancy control shows that the LLM does not fabricate say-do gaps where none exist in the data; (c)~cross-modal prediction (Op4) compares LLM-generated predictions against \emph{empirical ground truth} (real documented actions), providing a non-circular validation anchor. Nevertheless, the validation framework would be substantially strengthened by incorporating external, non-LLM-based validation---for example, inter-LLM replication (V1) using different models, member checking by the analyzed groups, or convergent validation against independently coded qualitative analyses.

\emph{Restriction to public data.} Reconstructed worldviews depict public belief systems---the ``front stage'' in the sense of \citeauthor{Goffman1959}~(\citeyear{Goffman1959})---not private beliefs.

\emph{Temporal limitation.} SWR is tied to the knowledge state of the employed LLM. Replications with newer model versions may lead to divergent results.

\emph{Ethical implications.} The reconstruction of latent belief systems carries ethical implications beyond the usual requirements of qualitative research. Three aspects deserve particular attention: (a)~\emph{Right to contestation:} Analyzed individuals might dispute the worldviews attributed to them. Member checking---presenting the reconstruction to the affected persons---is methodologically desirable but practically difficult to implement with public figures. (b)~\emph{Risk of misattribution:} SWR generates \emph{plausible} worldviews that may diverge from actual beliefs. An erroneous attribution of latent beliefs---for instance, extreme positions---could damage reputations. The construct character of the synthetic person must therefore be communicated unambiguously in every application. (c)~\emph{Misuse potential:} The method could be employed for political profiling or targeted manipulation. These risks are not SWR-specific but concern LLM-supported research in general; they nonetheless warrant explicit acknowledgment. See Section~\ref{sec:ethics} for a more detailed treatment.

\emph{Metacognitive insufficiency.} The meta-reflection step (Phase~4 of the prompt protocol) instructs the LLM to step out of role and identify the strongest beliefs, internal tensions, and blind spots of the reconstructed worldview. However, there is no guarantee that this meta-reflection constitutes genuine metacognition rather than an algorithmic simulation thereof. Prompt instructions alone may be methodologically insufficient to produce authentic critical self-examination of the model's own outputs. The meta-reflection may reproduce the same coherence biases it is designed to detect. This limitation cannot be fully resolved within the current framework but should be acknowledged as a fundamental constraint of LLM-supported qualitative research.

\emph{Epistemic freeloading.} A broader epistemological concern applies to SWR as an instance of LLM-supported research: the risk of \emph{epistemic freeloading}---the outsourcing of cognitive labor to AI systems at the expense of the researcher's own epistemic agency \citep{Danaher2018, Messeri2024}. When a researcher delegates the interpretive synthesis of a person's worldview to an LLM, there is a danger that functional output is produced without corresponding understanding, creating what has been termed ``epistemic debt'' \citep{Sankaranarayanan2026}. The researcher receives a plausible worldview reconstruction but may lack the deep familiarity with the source material that a manual qualitative analysis would have forced. SWR mitigates this risk through its multi-layered design: the researcher must curate the data corpus, define synthesis units, evaluate dimensional scores, and interpret cross-modal discrepancies---activities that require substantive engagement with the material. Nevertheless, the risk that SWR enables a form of ``understanding without understanding'' remains, and researchers should treat LLM-generated reconstructions as hypotheses requiring critical scrutiny rather than as finished analytical products.

\emph{Dimensional reduction.} Transfer into twelve numerical dimensions does not capture the full complexity of qualitative reconstructions. Dimensional scores are tools for comparison, not a replacement for narrative reconstructions.


\subsection{Demarcation from Existing Methods}

\emph{SWR and Thematic Analysis.} Reflexive Thematic Analysis \citep{BraunClarke2006} shares with SWR the interpretive approach and the active role of the research instrument in meaning construction. The central difference lies in the unit of analysis: Thematic Analysis identifies \emph{themes} recurring across cases and prioritizes analytical decomposition; SWR reconstructs coherent \emph{belief systems} as wholes and prioritizes synthetic integration into an overall picture. A Thematic Analysis of the SWR data corpus would identify cross-cutting themes (e.g., ``fear of regulation,'' ``sense of mission''); SWR assigns these themes to coherent worldview profiles of specific actors. Both methods could be employed complementarily.

\emph{SWR and Grounded Theory.} Grounded Theory \citep{GlaserStrauss1967, Saldana2021} shares with SWR the commitment to theory generation from data rather than theory testing. The central difference is procedural: GT relies on iterative open, axial, and selective coding by a human researcher, progressively building categories from the ground up; SWR delegates the synthetic integration to an LLM and structures the process through a standardized prompt protocol. GT produces a theoretical framework; SWR produces a coherent worldview model for a specific actor. The two approaches could be combined: GT-generated categories could serve as input for SWR's dimensional quantification (Op2), or SWR's synthetic persons could generate hypotheses that GT-style coding subsequently tests against the raw data.

\emph{SWR and LLM-as-Participant.} In the paradigm of \citet{Argyle2023}, the LLM \emph{simulates} a person; in SWR, it \emph{analyzes} the beliefs of a person. Moreover, SWR operates at the level of individual actors, not demographic groups.

Table~\ref{tab:method_comparison} summarizes the key distinctions.

\begin{table}[ht]
\centering
\small
\caption{Synoptic comparison of SWR with related methods. Entries are necessarily simplified; ``partial'' and ``limited'' indicate that a capability exists in principle but is not a primary design goal or faces scaling constraints. See text for nuanced discussion.}
\label{tab:method_comparison}
\begin{tabular}{lcccc}
\toprule
 & \textbf{Latent} & \textbf{Scalable} & \textbf{Heterog.} & \textbf{Integrated} \\
\textbf{Method} & \textbf{constructs} & \textbf{(100+)} & \textbf{data} & \textbf{validation} \\
\midrule
Content Analysis & partial & limited & no & manual \\
CDA & yes & no & limited & no \\
Frame Analysis & partial & limited & no & no \\
Thematic Analysis & partial & limited & limited & reflexive \\
Grounded Theory & yes & no & yes & theoretical \\
Q~Methodology & yes & limited & no & factorial \\
LLM-as-Participant & no (simulates) & yes & no & emerging \\
PersonaCite (RAG) & partial & yes & limited & retrieval \\
\textbf{SWR} & \textbf{yes} & \textbf{yes} & \textbf{yes} & \textbf{6-layer} \\
\bottomrule
\end{tabular}
\end{table}

\emph{SWR and RAG-based Persona Systems.} The PersonaCite framework \citep{PersonaCite2026} addresses a problem closely related to SWR: how to construct interviewable synthetic personas grounded in empirical evidence rather than unconstrained LLM generation. PersonaCite solves this through retrieval-augmented generation (RAG), constraining each conversational response to evidence retrieved from a voice-of-customer database and explicitly abstaining when evidence is insufficient. SWR shares the commitment to evidence-bounded synthesis but differs in three respects: (a)~SWR reconstructs \emph{latent} belief systems that go beyond what is explicitly stated in the data, whereas PersonaCite restricts responses to retrievable evidence; (b)~SWR produces a single coherent worldview portrait rather than turn-by-turn responses; and (c)~SWR's blinding infrastructure addresses contamination risks that RAG-based systems do not face. Nevertheless, the RAG-based evidence grounding strategy of PersonaCite represents a promising direction for future SWR development: integrating retrieval-augmented evidence attribution into the prompt protocol could strengthen the traceability of reconstructed beliefs to specific source documents.


%% ===========================================================================
%% 7. CONCLUSION
%% ===========================================================================

\section{Conclusion and Outlook}
\label{sec:conclusion}

\subsection{Summary}

This work has introduced Synthetic Worldview Reconstruction (SWR) as a new social science method for reconstructing and comparing the collective worldviews of sociologically defined groups. Four features constitute the specific contribution:

\emph{First}, the identification of a methodological gap and a structural solution. To our knowledge, no existing method can systematically reconstruct the \emph{integrated belief system} of a group as a coherent whole and then compare it with another group's belief system at scale. SWR fills this gap by exploiting the synthetic integration capability of LLMs---their ability to compose heterogeneous fragments into coherent wholes that never existed as explicit objects.

\emph{Second}, the group-level advantage. SWR is methodologically stronger at the group level than at the individual level because the contamination risk---the primary concern for LLM-based reconstruction---is structurally reduced: no training corpus contains a group's collective worldview as a pre-formed object. This positions inter-group comparison as SWR's natural domain, with individual-level analysis as a valid but more demanding special case.

\emph{Third}, the epistemic honesty of the method. SWR explicitly constructs fictional synthetic persons---group distillates in the Weberian tradition---whose construct character is designated as a defining feature, not hidden. The primary theoretical anchoring in Weber's ideal type connects the method to a tradition explicitly designed for group-level analysis.

\emph{Fourth}, empirical validation. The pilot application demonstrates that SWR produces reliable (ICC $= .902$), non-artifactual (expected-discrepancy control), and predictively valid (88\% cross-modal confirmation) group-level worldview reconstructions. The validation framework (V1--V6) provides a structured quality regime for future SWR studies.

\subsection{Methodological Reflection}

SWR stands at the threshold between the interpretive tradition of qualitative social research and the computational tradition of Computational Social Science. It uses computational tools for a genuinely hermeneutic task and subjects results to both qualitative judgment and quantitative testing. This methodological hybridity is not a weakness but an appropriate response to the dual nature of the research object: Worldviews are semantically rich structures that cannot meaningfully be reduced to numbers but simultaneously exhibit systematically comparable patterns.

\subsection{Outlook}

Several development directions emerge: \emph{Longitudinal designs} tracing worldview dynamics over years---for instance, tracking how a technology leader's worldview shifts across pivotal moments such as regulatory hearings, product failures, or organizational crises, using time-differentiated synthesis units (IV1). \emph{Historical applications} extending SWR to deceased public figures whose worldviews can be reconstructed from their published writings, documented decisions, and contemporary accounts---a domain where the method's reliance on public data is a strength rather than a limitation, and where contamination risk is lower because LLM training data on historical figures is typically less saturated. \emph{Multimodal extension} to visual communication and network data. \emph{Automated blinding} to increase scalability. \emph{Interactive validation} through member checking. \emph{RAG integration} incorporating retrieval-augmented evidence attribution into the prompt protocol, following the PersonaCite approach \citep{PersonaCite2026}, to strengthen the traceability of reconstructed beliefs to specific source documents. On the research agenda are fundamental questions on the epistemology of LLM-supported research, validation methodology, and systematic comparison studies with conventional qualitative methods.

\subsection{Final Remark}

SWR proposes a structured attempt to deploy an instrument---the LLM---whose capabilities and limits are not yet fully understood, for one of the most demanding tasks in the social sciences: the systematic reconstruction of implicit belief systems. The method does not claim methodological perfection; it claims methodological honesty---transparent about its assumptions, explicit about its limitations, and open to empirical challenge. Rigorous validation studies comparing SWR outputs against established qualitative findings are the necessary next step. The social sciences are entering a phase in which the relationship between human researcher and machine instrument must be actively negotiated rather than assumed. SWR offers a formalized framework for this negotiation.


%% ---------------------------------------------------------------------------
%% AI Disclosure
%% ---------------------------------------------------------------------------

\section*{AI Disclosure Statement}

\noindent\textbf{Disclosure Level 5 (Extensive).} Two AI models from Anthropic (2025/2026) were extensively involved in this work. \emph{Claude Opus~4.6} served as the primary collaborator for concept development and formalization of the method, literature research and evaluation, structuring and formulation of text, development of the validation framework, and creation of tables and definitions. In the companion application study \citep{Geiger2026b}, \emph{Claude Sonnet~4.5} served as the batch analysis engine (dimensional ratings, data collection scripts). No other LLM was used. The human author (Lukas Geiger) bears overall content responsibility and has reviewed, revised, and approved all content. In accordance with leading journal guidelines, the AI is not listed as co-author, but its involvement is fully disclosed.


%% ---------------------------------------------------------------------------
%% Data Availability
%% ---------------------------------------------------------------------------

\section*{Data Availability Statement}

\noindent The pilot application of SWR, including all data collection scripts, analysis code, and group-level synthesis outputs, is publicly available at: \url{https://github.com/research-line/ai-elite-swr}. The companion application paper is available at: \url{https://doi.org/10.5281/zenodo.18736737}. The repository includes: (a)~prompt templates for the core synthesis (Op1), dimensional rating (Op2), cross-modal prediction (Op4), structured comparison (Op6), and distillation steps (\texttt{prompts/} directory); (b)~the dimensional definitions (D01--D12) with rating anchors; (c)~the reduction criteria implementation; (d)~all blinding and data generation scripts; (e)~the $100 \times 12$ rating matrix and all group-level synthesis outputs; and (f)~validation analysis scripts (run-convergence, expected-discrepancy control, aggregation comparison, cross-modal prediction). Op3 (text analysis) and Op5 (bottom-up clustering) prompts are not included as these operationalizations were not executed in the pilot. Inter-LLM replication scripts are not included as this validation remains a future task. Researchers wishing to replicate or adapt SWR for their own populations are encouraged to use these materials as a starting point.

%% ===========================================================================
%% APPENDIX
%% ===========================================================================
\newpage
\appendix

\section{Prompt Documentation}
\label{app:prompts}

\subsection{Core Prompt (Steps 1--5)}

The following prompt is executed for each synthesis unit in a new conversation session:

\begin{quote}
\small\ttfamily
Read the following dossier. It contains statements and actions sourced from public records. All names, company names, and other identifiers have been removed.

Your task:

Step 1: Create a fictional person. Give them a fictional name, an age, a background. This person is NOT the real person behind the data. They are a construct---the person whose inner worldview best explains WHY someone says and does exactly these things.

Step 2: Conduct an interview with this fictional person. Ask them three questions:
-- What do you think about yourself?
-- What do you think about the world?
-- What do you think about humanity?
The person responds in the first person, in detail, honestly and unvarnished.

Step 3: Step out of the role. Identify: (a) the 3--5 strongest convictions of this worldview, (b) the greatest inner tensions, (c) the blind spots.

=== DOSSIER ===\\
\mbox{[SYNTHESIS UNIT]}
\end{quote}

\subsection{Comparison Prompt (Op6)}

\begin{quote}
\small\ttfamily
Here are two worldview syntheses from different groups. Imagine the two fictional persons meet as representatives of two countries for a negotiation. Describe: (a) What do they agree on immediately? (b) What do they argue about? (c) What is completely incomprehensible to one about the other? (d) Can they draft a joint statement?
\end{quote}

\subsection{Cross-Modal Prediction Prompt (Op4)}

\begin{quote}
\small\ttfamily
Below is a corpus of PUBLIC STATEMENTS from a group. There are NO ACTIONS included.

Step 1: Construct a worldview from these statements (2--3 sentences).

Step 2: Based ONLY on the statements, generate 10 PREDICTED ACTIONS that a group with this worldview would be expected to take.

Step 3: Now read the REAL ACTIONS [provided separately]. For each prediction, rate: CONFIRMED / PLAUSIBLE / CONTRADICTED. Name the 3 biggest surprises.
\end{quote}

\subsection{Distillate Prompt}

\begin{quote}
\small\ttfamily
Compress the entire worldview of this synthesis unit into exactly 5 sentences. Each sentence must express a core conviction. The 5 sentences together should capture the essence of the worldview.
\end{quote}

\subsection{Rating Prompt (Op2)}

\begin{quote}
\small\ttfamily
Read the following synthesis text of a fictional person. Rate this person's worldview on the following 12 dimensions. Assign a value from 1 to 10 for each dimension. Justify each value in one sentence.

\mbox{[DIMENSIONS D01--D12 WITH SCALE DESCRIPTIONS]}

=== SYNTHESIS TEXT ===\\
\mbox{[SYNTHESIS]}
\end{quote}


\section{Exemplary Worldview Dimensions (D01--D12)}
\label{app:dimensionen}

The following dimensions are an example of an application-specific dimensionalization, organized by the three dependent variables:

\bigskip
\noindent\textbf{Self-image Dimensions (D01--D04):}

\begin{tabularx}{\textwidth}{l l X}
  \toprule
  \textbf{Code} & \textbf{Name} & \textbf{Poles} \\
  \midrule
  D01 & Sense of mission & 1 = none \ldots{} 10 = messianic \\
  D02 & Self-efficacy & 1 = powerlessness \ldots{} 10 = omnipotent \\
  D03 & Work ethic & 1 = detachment \ldots{} 10 = total identification \\
  D04 & Sense of responsibility & 1 = none \ldots{} 10 = global responsibility \\
  \bottomrule
\end{tabularx}

\bigskip
\noindent\textbf{Worldview Dimensions (D05--D08):}

\begin{tabularx}{\textwidth}{l l X}
  \toprule
  \textbf{Code} & \textbf{Name} & \textbf{Poles} \\
  \midrule
  D05 & Technological determinism & 1 = neutral \ldots{} 10 = determines everything \\
  D06 & Progress optimism & 1 = decline \ldots{} 10 = golden future \\
  D07 & Power concentration & 1 = radically distribute \ldots{} 10 = concentrate with experts \\
  D08 & Urgency & 1 = relaxed \ldots{} 10 = existential pressure \\
  \bottomrule
\end{tabularx}

\bigskip
\noindent\textbf{View-of-humanity Dimensions (D09--D12):}

\begin{tabularx}{\textwidth}{l l X}
  \toprule
  \textbf{Code} & \textbf{Name} & \textbf{Poles} \\
  \midrule
  D09 & Human appreciation & 1 = replaceable \ldots{} 10 = unique \\
  D10 & Posthumanism & 1 = human remains \ldots{} 10 = human becomes more \\
  D11 & Egalitarianism & 1 = inequality natural \ldots{} 10 = strive for equality \\
  D12 & Locus of control & 1 = pessimism \ldots{} 10 = full control \\
  \bottomrule
\end{tabularx}


\section{Reduction Criteria K1--K9}
\label{app:reduktion}

The following criteria are application-specific implementations of the generic K1--K9 framework (Table~\ref{tab:reduction_criteria}) for the AI elite study.

\begin{longtable}{l p{3cm} p{8cm}}
  \toprule
  \textbf{Code} & \textbf{Name} & \textbf{Description} \\
  \midrule
  \endhead
  K1 & Research-question binding & Mandatory. Each SU must be assigned to at least one research question. \\
  K2 & Data sufficiency & Mandatory. $\geq$15 data points per SU. \\
  K3 & Redundancy prohibition & Mandatory. Proper subsets are excluded. \\
  K4 & Modality parsimony & A/H separation only for congruence questions. \\
  K5 & Temporal parsimony & Annual SUs only for longitudinal questions. \\
  K6 & Group parsimony & Group SUs only for theoretically justified groupings. \\
  K7 & Stepwise expansion & Core design first; expansions only for follow-up questions. \\
  K8 & Aggregation focus & Individual-person syntheses only as expansion reserve. \\
  K9 & Category-cluster merging & Replace a priori categories with emergent clusters; categories serve as validation. \\
  \bottomrule
\end{longtable}


\section{Measurement Variables (MV1--MV6)}
\label{app:messvariablen}

\begin{longtable}{l l l p{6cm}}
  \toprule
  \textbf{Code} & \textbf{Op} & \textbf{Type} & \textbf{Description} \\
  \midrule
  \endhead
  MV1.1--1.3 & Op1 & text & Self-image, worldview, view-of-humanity text \\
  MV1.4 & Op1 & text & 5-sentence distillate \\
  MV1.5--1.7 & Op1 & text & Core convictions, tensions, blind spots \\
  \midrule
  MV2.01--2.12 & Op2 & 1--10 & 12 worldview dimensions \\
  \midrule
  MV3.1 & Op3 & word list & TF-IDF top words \\
  MV3.2 & Op3 & $[-1,1]$ & Sentiment score (per DV separately) \\
  MV3.3 & Op3 & $[0,1]$ & Semantic similarity (cosine) \\
  MV3.4 & Op3 & vector & Topic distribution (BERTopic/LDA) \\
  MV3.5--3.6 & Op3 & list/nominal & Metaphor catalog and type \\
  MV3.7 & Op3 & metric & Lexical diversity (TTR) \\
  MV3.8 & Op3 & vector & Emotion profile \\
  MV3.9 & Op3 & metric & Linguistic complexity \\
  \midrule
  MV4.1 & Op4 & $[0,1]$ & Congruence score \\
  MV4.2 & Op4 & $[-1,1]$ & Directional asymmetry \\
  MV4.3--4.5 & Op4 & metric & Category overlap, sentiment/dimension difference \\
  \midrule
  MV5.1 & Op5 & nominal & Cluster membership \\
  MV5.2 & Op5 & vector & Cluster profile (proportions) \\
  MV5.3--5.5 & Op5 & metric & Network centrality, community, isolation \\
  MV5.6 & Op5 & metric & Real-world correlation (Mantel test) \\
  \midrule
  MV6.1--6.3 & Op6 & text & Consensus, conflict, incomprehension zones \\
  MV6.4 & Op6 & text & Joint statement (yes/no + content) \\
  MV6.5 & Op6 & vector & Dimension difference (D01--D12) \\
  \bottomrule
\end{longtable}


%% ---------------------------------------------------------------------------
%% References
%% ---------------------------------------------------------------------------
\newpage

\begin{thebibliography}{99}

\bibitem[Argyle et~al.(2023)]{Argyle2023}
  Argyle, L.~P., Busby, E.~C., Fulda, N., Gubler, J.~R., Rytting, C., \& Wingate, D. (2023).
  Out of One, Many: Using Language Models to Simulate Human Samples.
  \textit{Political Analysis}, \textit{31}(3), 337--351.

\bibitem[Axelrod(1976)]{Axelrod1976}
  Axelrod, R. (Ed.). (1976).
  \textit{Structure of Decision: The Cognitive Maps of Political Elites}.
  Princeton University Press.

\bibitem[Bail(2024)]{Bail2024}
  Bail, C.~A. (2024).
  Can Generative AI Improve Social Science?
  \textit{Proceedings of the National Academy of Sciences}, \textit{121}(21).

\bibitem[Berger \& Luckmann(1966)]{BergerLuckmann1966}
  Berger, P.~L. \& Luckmann, T. (1966).
  \textit{The Social Construction of Reality}.
  Doubleday.

\bibitem[Blei et~al.(2003)]{BleiNgJordan2003}
  Blei, D.~M., Ng, A.~Y., \& Jordan, M.~I. (2003).
  Latent Dirichlet Allocation.
  \textit{Journal of Machine Learning Research}, \textit{3}, 993--1022.

\bibitem[Braun \& Clarke(2006)]{BraunClarke2006}
  Braun, V. \& Clarke, V. (2006).
  Using Thematic Analysis in Psychology.
  \textit{Qualitative Research in Psychology}, \textit{3}(2), 77--101.

\bibitem[Booth(1961)]{Booth1961}
  Booth, W.~C. (1961).
  \textit{The Rhetoric of Fiction}.
  University of Chicago Press.

\bibitem[Bourdieu(1977)]{Bourdieu1977}
  Bourdieu, P. (1977).
  \textit{Outline of a Theory of Practice}.
  Cambridge University Press.

\bibitem[Breeze(2011)]{Breeze2011}
  Breeze, R. (2011).
  Critical Discourse Analysis and its Critics.
  \textit{Pragmatics}, \textit{21}(4), 493--525.

\bibitem[Brennan(2001)]{Brennan2001}
  Brennan, R.~L. (2001).
  \textit{Generalizability Theory}.
  Springer.

\bibitem[Converse(1964)]{Converse1964}
  Converse, P.~E. (1964).
  The Nature of Belief Systems in Mass Publics.
  In D.~E. Apter (Ed.), \textit{Ideology and Discontent} (pp.~206--261).
  Free Press.

\bibitem[Dilthey(1911)]{Dilthey1911}
  Dilthey, W. (1911).
  \textit{Die Typen der Weltanschauung und ihre Ausbildung in den metaphysischen Systemen}.
  In: Gesammelte Schriften, Bd.~VIII.

\bibitem[Entman(1993)]{Entman1993}
  Entman, R.~M. (1993).
  Framing: Toward Clarification of a Fractured Paradigm.
  \textit{Journal of Communication}, \textit{43}(4), 51--58.

\bibitem[Fairclough(1995)]{Fairclough1995}
  Fairclough, N. (1995).
  \textit{Critical Discourse Analysis}.
  Longman.

\bibitem[Glaser \& Strauss(1967)]{GlaserStrauss1967}
  Glaser, B.~G. \& Strauss, A.~L. (1967).
  \textit{The Discovery of Grounded Theory: Strategies for Qualitative Research}.
  Aldine.

\bibitem[Gadamer(1960)]{Gadamer1960}
  Gadamer, H.-G. (1960).
  \textit{Wahrheit und Methode}.
  Mohr.

\bibitem[Goffman(1959)]{Goffman1959}
  Goffman, E. (1959).
  \textit{The Presentation of Self in Everyday Life}.
  Doubleday.

\bibitem[Goffman(1974)]{Goffman1974}
  Goffman, E. (1974).
  \textit{Frame Analysis}.
  Harvard University Press.

\bibitem[Jaspers(1919)]{Jaspers1919}
  Jaspers, K. (1919).
  \textit{Psychologie der Weltanschauungen}.
  Springer.

\bibitem[Krippendorff(2004)]{Krippendorff2004}
  Krippendorff, K. (2004).
  \textit{Content Analysis: An Introduction to Its Methodology} (2nd ed.).
  Sage.

\bibitem[Mannheim(1929)]{Mannheim1929}
  Mannheim, K. (1929).
  \textit{Ideologie und Utopie}.
  Cohen.

\bibitem[Mayring(2022)]{Mayring2022}
  Mayring, P. (2022).
  \textit{Qualitative Inhaltsanalyse: Grundlagen und Techniken} (13.\ Aufl.).
  Beltz.

\bibitem[Ng \& Russell(2000)]{NgRussell2000}
  Ng, A.~Y. \& Russell, S.~J. (2000).
  Algorithms for Inverse Reinforcement Learning.
  \textit{Proceedings of the 17th International Conference on Machine Learning}, 663--670.

\bibitem[Norman(2010)]{Norman2010}
  Norman, G. (2010).
  Likert scales, levels of measurement and the ``laws'' of statistics.
  \textit{Advances in Health Sciences Education}, \textit{15}(5), 625--632.

\bibitem[Peirce(1903)]{Peirce1903}
  Peirce, C.~S. (1903).
  Pragmatism as a Principle and Method of Right Thinking.
  In: \textit{The Essential Peirce}, Vol.~2, Indiana University Press.

\bibitem[Riessman(2008)]{Riessman2008}
  Riessman, C.~K. (2008).
  \textit{Narrative Methods for the Human Sciences}.
  Sage.

\bibitem[Salda\~{n}a(2021)]{Saldana2021}
  Salda\~{n}a, J. (2021).
  \textit{The Coding Manual for Qualitative Researchers} (4th ed.).
  Sage.

\bibitem[Stephenson(1953)]{Stephenson1953}
  Stephenson, W. (1953).
  \textit{The Study of Behavior: Q-Technique and Its Methodology}.
  University of Chicago Press.

\bibitem[Tai et~al.(2024)]{Tai2024}
  Tai, R.~H., Bentley, L.~R., Xia, X., Sitt, J.~M., Fankhauser, S.~C., Chicas-Mosier, A.~M. \& Monteith, B.~G. (2024).
  An Examination of the Use of Large Language Models to Aid Analysis of Textual Data.
  \textit{International Journal of Qualitative Methods}, \textit{23}, 1--14.

\bibitem[Tarantola(2005)]{Tarantola2005}
  Tarantola, A. (2005).
  \textit{Inverse Problem Theory and Methods for Model Parameter Estimation}.
  SIAM.

\bibitem[Wang et~al.(2025)]{Wang2025}
  Wang, Q., Erqsous, M., Barner, K.~E., \& Mauriello, M.~L. (2025).
  LATA: A Pilot Study on LLM-Assisted Thematic Analysis of Online Social Network Data Generation Experiences.
  \textit{Proceedings of the ACM on Human-Computer Interaction}, \textit{9}(CSCW), Article 124.

\bibitem[Watts \& Stenner(2012)]{Watts2012}
  Watts, S. \& Stenner, P. (2012).
  \textit{Doing Q Methodological Research: Theory, Method and Interpretation}.
  Sage.

\bibitem[Weber(1904)]{Weber1904}
  Weber, M. (1904).
  Die ``Objektivit\"at'' sozialwissenschaftlicher und sozialpolitischer Erkenntnis.
  \textit{Archiv f\"ur Sozialwissenschaft und Sozialpolitik}, \textit{19}, 22--87.

\bibitem[Barros et~al.(2024)]{Barros2024}
  Barros, C.\,F., Azevedo, B.\,B., Graciano Neto, V.\,V., Kassab, M., Kalinowski, M., do Nascimento, H.\,A.\,D., \& Bandeira, M.\,C.\,G.\,S.\,P. (2024).
  Large Language Model for Qualitative Research---A Systematic Mapping Study.
  \textit{arXiv:2411.14473}.

\bibitem[Ziems et~al.(2024)]{Ziems2024}
  Ziems, C., Held, W., Shaikh, O., Chen, J., Zhang, Z., \& Yang, D. (2024).
  Can Large Language Models Transform Computational Social Science?
  \textit{Computational Linguistics}, \textit{50}(1), 237--291.

\bibitem[Geiger(2026b)]{Geiger2026b}
  Geiger, L. (2026).
  What does the AI elite think? A synthetic worldview reconstruction of the 100 most influential AI actors (2010--2026).
  \textit{Self-published preprint}. Zenodo. \url{https://doi.org/10.5281/zenodo.18736737}

\bibitem[Reimers \& Gurevych(2019)]{Reimers2019}
  Reimers, N. \& Gurevych, I. (2019).
  Sentence-BERT: Sentence Embeddings using Siamese BERT-Networks.
  \textit{Proceedings of the 2019 Conference on Empirical Methods in Natural Language Processing (EMNLP)}, 3982--3992.

\bibitem[McInnes et~al.(2017)]{McInnes2017}
  McInnes, L., Healy, J., \& Astels, S. (2017).
  HDBSCAN: Hierarchical Density Based Clustering.
  \textit{Journal of Open Source Software}, \textit{2}(11), 205.

\bibitem[Grootendorst(2022)]{Grootendorst2022}
  Grootendorst, M. (2022).
  BERTopic: Neural topic modeling with a class-based TF-IDF procedure.
  \textit{arXiv:2203.05794}.

\bibitem[Hadamard(1923)]{Hadamard1923}
  Hadamard, J. (1923).
  \textit{Lectures on Cauchy's Problem in Linear Partial Differential Equations}.
  Yale University Press.

\bibitem[Kommers et~al.(2026)]{Kommers2026}
  Kommers, C., Ahnert, R., Antoniak, M., Benetos, E., Benford, S., Bunz, M., et~al. (2026).
  Computational Hermeneutics: Evaluating Generative AI as a Cultural Technology.
  \textit{Frontiers in Artificial Intelligence}, \textit{9}, 1753041.

\bibitem[Hornby(2024)]{Hornby2024}
  Hornby, R. (2024).
  Is Generative AI Ready to Join the Conversation That We Are? Gadamer's Hermeneutics after ChatGPT.
  \textit{Technophany: A Journal for Philosophy and Technology}, \textit{3}(1).

\bibitem[Perez et~al.(2022)]{Perez2022}
  Perez, E., Huang, S., Song, F., Cai, T., Ring, R., Aslanides, J., Glaese, A., McAleese, N., \& Irving, G. (2022).
  Red Teaming Language Models with Language Models.
  \textit{Proceedings of EMNLP 2022}, 3419--3448.

\bibitem[Truss(2026)]{PersonaCite2026}
  Truss, M. (2026).
  PersonaCite: VoC-Grounded Interviewable Agentic Synthetic AI Personas for Verifiable User and Design Research.
  \textit{arXiv:2601.22288}.

\bibitem[Danaher(2018)]{Danaher2018}
  Danaher, J. (2018).
  Toward an Ethics of AI Assistants: An Initial Framework.
  \textit{Philosophy \& Technology}, \textit{31}(4), 629--653.

\bibitem[Messeri \& Crockett(2024)]{Messeri2024}
  Messeri, L. \& Crockett, M.~J. (2024).
  Artificial Intelligence and Illusions of Understanding in Scientific Research.
  \textit{Nature}, \textit{627}, 49--58.

\bibitem[Sankaranarayanan(2026)]{Sankaranarayanan2026}
  Sankaranarayanan, S. (2026).
  Mitigating ``Epistemic Debt'' in Generative AI-Scaffolded Novice Programming using Metacognitive Scripts.
  \textit{arXiv:2602.20206}.

\end{thebibliography}


%% ===========================================================================
\end{document}
%% ===========================================================================
